\documentclass[8pt,letterpaper]{article}
\usepackage[margin=0.25in,top=0.2in,bottom=0.15in]{geometry}
\usepackage{multicol}
\usepackage[most]{tcolorbox}
\usepackage{xcolor}
\usepackage{helvet}
\usepackage{hyperref}
\usepackage{fontawesome5}
\usepackage[nolinks]{qrcode}
\renewcommand{\familydefault}{\sfdefault}
\setlength{\parindent}{0pt}
\setlength{\parskip}{0pt}
\pagestyle{empty}
\hypersetup{colorlinks=true,urlcolor=blue!70!black,linkcolor=black}
\renewcommand{\UrlBreaks}{\do\/\do-\do.\do_\do=}

% ── Colors ───────────────────────────────────────────────────────────────────
\definecolor{headerblue}{RGB}{52,79,115}
\definecolor{codebg}{RGB}{245,246,248}
\definecolor{codetext}{RGB}{30,30,30}
\definecolor{codeborder}{RGB}{200,205,215}
\definecolor{discordpurple}{RGB}{88,101,242}

% ── Verbatim code listing inside tcolorbox ───────────────────────────────────
\tcbuselibrary{listings,skins}

\newtcblisting{codeblock}{
  colback=codebg, colframe=codeborder,
  arc=2pt, boxrule=0.4pt,
  left=2pt, right=2pt, top=1pt, bottom=1pt,
  listing only,
  listing options={
    basicstyle=\fontsize{5.8}{7.0}\ttfamily\color{codetext},
    breaklines=true,
    keepspaces=true,
    columns=fullflexible,
  },
  before=\vspace{1pt},
  after=\vspace{0pt},
}

% ── Section box ──────────────────────────────────────────────────────────────
\newtcolorbox{secbox}[1]{
  colback=white, colframe=headerblue,
  title={\small\bfseries\color{white}#1},
  colbacktitle=headerblue,
  boxrule=0.5pt,
  arc=3pt,
  left=2pt, right=2pt, top=0pt, bottom=1pt,
  toptitle=1pt, bottomtitle=1pt,
  before=\vspace{1pt},
  after=\vspace{1pt},
}

% ─────────────────────────────────────────────────────────────────────────────
\begin{document}
% ── Title ────────────────────────────────────────────────────────────────────
{\LARGE\bfseries Secure DevOps for Ansible%
  \hfill\normalsize\color{gray}One-Page Reference}\\[-3pt]
\textcolor{headerblue}{\rule{\linewidth}{1.5pt}}\\[1pt]
{\small
  Flow:\enspace SSH key $\to$ SSH-Agent $\to$ GPG key\,/\,GPG-Agent $\to$
  \texttt{pass} $\to$ vault helper $\to$ \texttt{ansible.cfg} $\to$
  encrypted \texttt{vault.yml} $\to$ \texttt{ansible-playbook}}\\[2pt]

% ── Two columns ──────────────────────────────────────────────────────────────
\setlength{\columnsep}{8pt}
\begin{multicols}{2}

\begin{secbox}{1) Generate SSH and GPG Keys}
\begin{codeblock}
# SSH key for remote access
ssh-keygen -t ed25519 -C "ansible-demo" \
  -f ~/.ssh/id_ed25519

# GPG key for pass / vault password storage
gpg --full-generate-key
\end{codeblock}
\end{secbox}

\begin{secbox}{2) Start Agents and Load Keys}
\begin{codeblock}
# SSH agent — add to ~/.bashrc for persistence
eval "$(ssh-agent -s)"
ssh-add ~/.ssh/id_ed25519

# GPG agent — ensure it is running
gpg-connect-agent /bye
\end{codeblock}
\end{secbox}

\begin{secbox}{3) \textasciitilde/.ssh/config}
\begin{codeblock}
Host ansible
    HostName 192.168.1.10
    User ansible_user
    IdentityFile ~/.ssh/id_ed25519
    AddKeysToAgent yes

Host server1
    HostName 192.168.1.101
    User ansible_user
    IdentityFile ~/.ssh/id_ed25519
    AddKeysToAgent yes
    StrictHostKeyChecking accept-new
\end{codeblock}
\end{secbox}

\begin{secbox}{4) pass --- The Unix Password Manager}
\begin{codeblock}
# Find your GPG key ID
gpg --list-secret-keys --keyid-format LONG

# Initialize pass with your key ID
pass init YOURKEYID

# Store and retrieve the vault password
pass insert ansible/vault_password
pass ansible/vault_password
\end{codeblock}
\end{secbox}

\begin{secbox}{5) Vault Password Helper Script}
\begin{codeblock}
mkdir -p ~/bin
cat > ~/bin/get_vault_pass.sh <<'EOF'
#!/bin/sh
pass ansible/vault_password
EOF
chmod 700 ~/bin/get_vault_pass.sh
\end{codeblock}
\end{secbox}

\begin{secbox}{6) Configure Ansible (\texttt{ansible.cfg})}
\begin{codeblock}
[defaults]
inventory           = inventory.ini
interpreter_python  = /usr/bin/python3
vault_password_file = ~/bin/get_vault_pass.sh

[ssh_connection]
ssh_args = -o StrictHostKeyChecking=accept-new
\end{codeblock}
\end{secbox}

\begin{secbox}{7) Encrypted Secrets in \texttt{vault.yml}}
\begin{codeblock}
# Create / edit / view the encrypted vault
ansible-vault create vault.yml
ansible-vault edit   vault.yml
ansible-vault view   vault.yml

# Plaintext values inside vault.yml:
ansible_become_password: YourSudoPass
ansible_user_password:   YourUserPass

# On disk it is AES-256 encrypted:
$ANSIBLE_VAULT;1.1;AES256
6136616539613930...
\end{codeblock}
\end{secbox}

\begin{secbox}{8) Playbook Uses \texttt{vault.yml}}
\begin{codeblock}
---
- name: Update packages securely
  hosts: servers     # remote_user via ~/.ssh/config
  become: yes
  vars_files:
    - vault.yml
  tasks:
    - name: Update apt cache
      apt:
        update_cache: yes
\end{codeblock}
\end{secbox}

\begin{secbox}{9) Run, Verify, and Harden}
\begin{codeblock}
ansible servers -m ping
ansible-playbook update.yml

# /etc/ssh/sshd_config hardening
PermitRootLogin      prohibit-password
PubkeyAuthentication yes
PasswordAuthentication no
\end{codeblock}
\end{secbox}

\end{multicols}

\vspace{1pt}
\textcolor{headerblue}{\rule{\linewidth}{0.8pt}}
\vspace{2pt}

% ── References ───────────────────────────────────────────────────────────────
{\fontsize{6.5}{8.5}\selectfont
\begin{tabular}{@{}p{0.245\linewidth}p{0.245\linewidth}p{0.245\linewidth}p{0.245\linewidth}@{}}
\textbf{\color{headerblue}Ansible} &
\textbf{\color{headerblue}SSH} &
\textbf{\color{headerblue}GPG / pass} &
\textbf{\color{headerblue}Hardening} \\[1pt]
\textbf{Ansible for DevOps} — Geerling\newline
{\fontsize{5.5}{7}\selectfont\href{https://leanpub.com/ansible-for-devops}{leanpub.com/ansible-for-devops}} &
\textbf{SSH Mastery} — M.W. Lucas\newline
{\fontsize{5.5}{7}\selectfont\href{https://mwl.io/nonfiction/tools\#ssh}{mwl.io/nonfiction/tools\#ssh}} &
\textbf{pass} — Unix Password Manager\newline
{\fontsize{5.5}{7}\selectfont\href{https://passwordstore.org}{passwordstore.org}} &
\textbf{CIS Benchmarks}\newline
{\fontsize{5.5}{7}\selectfont\href{https://cisecurity.org/cis-benchmarks}{cisecurity.org/cis-benchmarks}} \\[3pt]
\textbf{Ansible Vault Guide}\newline
{\fontsize{5.5}{7}\selectfont\href{https://docs.ansible.com/ansible/latest/vault_guide}{docs.ansible.com/ansible/latest/vault\_guide}} &
\textbf{Mozilla OpenSSH Guidelines}\newline
{\fontsize{5.5}{7}\selectfont\href{https://wiki.mozilla.org/Security/Guidelines/OpenSSH}{wiki.mozilla.org/Security/Guidelines/OpenSSH}} &
\textbf{OpenPGP Best Practices}\newline
{\fontsize{5.5}{7}\selectfont\href{https://riseup.net/en/security/message-security/openpgp/best-practices}{riseup.net/en/security/message-security/openpgp/best-practices}} &
\textbf{NIST SP 800-123}\newline
{\fontsize{5.5}{7}\selectfont\href{https://csrc.nist.gov/publications/detail/sp/800-123/final}{csrc.nist.gov/publications/detail/sp/800-123/final}} \\
\end{tabular}
}

\vspace{2pt}
{\small\color{gray}Focused on SSH/GPG keys, agents, \texttt{pass}, and Ansible Vault --- no plaintext passwords anywhere in the flow.}\\[1pt]
\textcolor{headerblue}{\rule{\linewidth}{0.8pt}}
\vspace{3pt}

% ── Footer with Discord QR code ──────────────────────────────────────────────
\begin{minipage}[c]{0.72\linewidth}
  {\normalsize\textbf{\textcolor{discordpurple}{\faDiscord\enspace TCC Cybersecurity Club}}}\\[1pt]
  {\small Join us on Discord:\enspace
    \href{https://discord.gg/k3pgCGX}%
    {\textcolor{discordpurple}{\textbf{\texttt{discord.gg/k3pgCGX}}}}}\\[4pt]
  {\fontsize{6.5}{8}\selectfont\color{gray}
    \textcopyright\ 2026 Joel Kirch \quad
    \href{https://creativecommons.org/licenses/by/4.0/}%
    {\faCreativeCommons\ \faCreativeCommonsBy\ CC BY 4.0}}
\end{minipage}%
\hfill%
\begin{minipage}[c]{0.22\linewidth}\raggedleft
  \qrcode[height=1.5cm]{https://discord.gg/k3pgCGX}
\end{minipage}

\end{document}
